\section{Discussion: Deployment, Extensions, and Expected Findings}
\label{sec:discussion}

\paragraph{A staged path to deployment.}
The proposed method can be developed in stages that make its intermediate behavior inspectable. The existing demonstration and episode memory banks supply source records, while the execution queue supplies request identities and receipts. The first stage adds phase annotation, evidence references, graph validation, and deterministic serialization. Offline prefix replay can then test whether the recorded conditions were observable when a decision was made. The second stage adds graph retrieval and adaptation, evaluated first as plan generation with the controller held fixed. The third introduces delegated execution, using recorded traces to check supervisor authority before robot trials. Finally, backend adapters expose the common contract for the cross-platform comparison. Each stage has a measurable output: valid records, applicable retrievals, justified edits, correctly authorized actions, or conforming interfaces.

This sequence also identifies where engineering effort may dominate. Phase boundaries and predicates must be detailed enough to distinguish decisions but stable enough to annotate consistently. An overly fine vocabulary increases annotation and retrieval cost; an overly coarse one can merge a failed action with the inspection that explains its recovery. A useful initial deployment can therefore restrict the skill vocabulary and object roles while retaining explicit uncertainty. Expanding coverage should follow demonstrated annotation agreement and reliable predicate evaluation. The cost of creating and maintaining the bank belongs in the assessment of practical value, alongside inference and execution cost.

\paragraph{Reusable procedures with inspectable evidence.}
The main deployment value would be the ability to inspect why a prior interaction applies to a new decision. A retrieved recovery includes its triggering failure, the observation that followed, and the conditions required for another attempt. A plan edit records which source supported it, while the execution trace records whether its anticipated effect occurred. This chain could support diagnosis, correction, and reuse across tasks without rereading entire episodes. For example, a failed drawer-transfer attempt can be traced to a missing opening prerequisite, a wrong object binding, or a false completion judgment. Trace readability and attribution correctness would measure this diagnostic value alongside task performance.

\paragraph{Growing the experience library.}
Longer deployments raise a selection problem as well as a storage problem. Repeated successful episodes may add little coverage, whereas an uncommon failure can reveal a missing condition or recovery branch. A future admission policy could prioritize new role configurations, previously unsupported transitions, and corrected predicate evaluations. Duplicate detection and source-level versioning would keep similar traces from dominating retrieval. Contradictory outcomes should retain their conditions and provenance until the difference can be explained. The chronological evaluation can measure whether these admission policies add coverage while retaining competence on earlier tasks.

\paragraph{Uncertain branches and multiple resources.}
The current three-valued conditions provide a direct response to incomplete observations: acquire permitted evidence or request revision. A useful extension is to choose observations according to which unresolved branch they can distinguish, accounting for their time and motion cost. This requires calibrated evaluation of the observations and should be tested against the simpler fixed inspection policy. Multi-arm operation is another scoped extension. Parallel graph phases would need explicit ownership of arms, grippers, shared workspaces, and transferred objects. Resource contracts could state when these become available, while local controllers handle synchronized motion. Deploying these phases requires backend support for coordinated motion and contact control.

\paragraph{Findings that would support the method.}
The hoped-for finding is that preserving phase conditions and outcome evidence helps robots reuse procedures across task compositions. The matched sequence comparison and retrieval-by-execution experiment would locate this benefit in retrieval, adaptation, or execution. Equal sequence performance would weaken the case for explicit graph enforcement; no gain from structural edits on held-out compositions would weaken the case for adaptation. These outcomes would favor a simpler representation or adaptation procedure. Delegation would have practical value if it reduced total cost while maintaining independently measured task success and completion reliability. The systems study would establish how much application code can be reused, how reliably adapters preserve the contract, and whether MCP's integration benefits justify its overhead. Such findings would guide investment in the graph bank, execution hierarchy, and backend support according to their demonstrated contribution to reliable manipulation.

% CHAPTER_SUMMARY: Deployment proceeds from existing memory and queue infrastructure through annotation, graph retrieval/adaptation, delegated execution, and backend adapters. Practical value comes from reusable procedures whose applicability and observed effects can be inspected. Scoped extensions concern library admission, uncertainty-driven observation, and multi-arm resource contracts. Desired findings are conditional: matched representation tests must establish structural value, held-out compositions must establish adaptation value, delegation must maintain independent success at lower cost, and systems tests must separate semantic reuse from MCP effects. No empirical conclusion is asserted.
